\section{Discussion}
\label{sec:discussion}

\paragraph{Comparability across phases.}
The three phases differ in evaluation conditions --- Phase~1 is zero-shot on the
full test split (8,231 items), Phase~2 trains to \texttt{max\_steps}~$=500$ and
evaluates on three datasets $\times$ 3 seeds over the full test split, and
Phase~3 trains to \texttt{max\_steps}~$=200$ and evaluates on the PathVQA
validation split, up to 500 samples. Direct comparison of absolute accuracy
across phases is therefore invalid: the difference between the fine-tuned
accuracy in Phase~2 (Qwen3-VL-2B, 0.5845) and the best \texttt{val\_accuracy} in
Phase~3 (0.4700) does not represent a performance drop. The discussion below
cross-references only rankings, patterns, and directions of gaps observed within
each phase.

\paragraph{Resource consumption and performance are not monotonically related.}
The best-performing model on pooled zero-shot accuracy, Qwen3-VL-2B (0.3843,
2B active / 2B total), simultaneously uses the least memory ($\sim$4,500\,MB) and
responds fastest, while Gemma4-E2B (0.1708, 2.3B active / 5.1B total), which uses
the most memory ($\sim$13,930\,MB), ranks last with a statistically pronounced gap
from the other three. At this scale, architecture and pretraining configuration
dominate medical VQA performance more than active parameter count or memory
consumption, and a choice satisfying both performance and efficiency exists for
resource-constrained environments. The axis of interpretation must, however, be
made explicit: only Gemma4-E2B has a mixture-of-experts structure in which active
parameters differ from stored parameters, and our selection criterion for
``lightweight'' took active parameters as its basis. Measured by stored
parameters Gemma4-E2B is the largest model, so its last place may equally be read
as ``the model with the largest stored scale ranked last'' rather than ``a small
model beat a large one.'' The claim concerns performance relative to
inference-time resource consumption and must not be extended to a general
relationship between parameter scale as such and performance. Because the ranking
survived removal of samples suspected of pretraining exposure, it is also not
explained by data contamination.

\paragraph{Heterogeneity of fine-tuning effects, and the primacy of model
selection.}
The proposition that fine-tuning improves performance does not hold independently
of the model, and its direction coincides with the zero-shot ranking: the models
ranked first and second zero-shot improved significantly ($d = +1.620$ and
$+2.646$), whereas the third and fourth either deteriorated significantly
($d = -2.284$) or failed to improve ($d = -0.652$, n.s.). Fine-tuning did not
function as a means of lifting a weak base model, which suggests in practical
terms that base-model selection is a decision preceding fine-tuning design; with
only four models covered, whether this correspondence is a general law remains
unconfirmed. Overlaying the forgetting analysis reveals a trade-off: the Qwen
family gave up 4--7\% of general capability on VQAv2, SmolVLM2-2.2B almost
entirely preserved general capability while deteriorating in the domain, and
Gemma4-E2B suffered a double loss, losing 51.5\% of general capability without
any domain gain. Domain specialization and preservation of general capability
were not achieved simultaneously within the scope of this experiment. This is an
observed correlation; no causal relationship was verified.

\paragraph{The closed--open gap is unresolved across all three phases.}
In Phase~1, PathVQA open accuracy was 0.027--0.061 across all models against
closed accuracy of 0.45--0.79, with all four models scoring 0.0000 on the
diagnosis and temporal types. In Phase~3, even the best configuration among 810
trials reached only \texttt{val\_open\_acc} 0.1562 against
\texttt{val\_closed\_acc} 0.8156 (the best configuration of the four-strategy
comparison gives 0.8115 / 0.1445). Absolute values across phases are not
comparable, yet the ratio relationship --- open below 20\% of closed --- holds in
every phase (about 9.5\% for Qwen3-VL-2B on PathVQA in Phase~1, about
17.8--19.2\% for the best Phase~3 configurations). That the ratio persisted even
in the re-experiment reinforces that this gap is not resolvable at the level of
hyperparameter search. It matters because the clinically important question types
are concentrated in the open-ended category. A type-level paired decomposition on
an identical 500-item sample (Appendix~\ref{app:tables},
Table~\ref{tab:appgain}) confirms the concentration: the two lowest-weighted
types, yes\_no (0.5) and description (0.6), account for 84.6\% of the total gain
and make up 91\% of the evaluation sample. Three qualifications must be read
alongside it. The sample sizes of the highest-weighted types fall short of what a
judgement requires --- of the 500 items, diagnosis has 3, measurement has 4, and
temporal none --- so their zero contribution means less that they failed to
improve than that the sample cannot decide whether they did. The per-type gain is
in fact largest among the higher-weighted types: location, weighted 0.8, rose
from 0.2381 to 0.7143, the largest gain of any type, its 15.4\% share reflecting
its 21-item sample rather than its performance. And weighting by clinical
importance does not cancel the gain: WCA rose from 0.1527 zero-shot to 0.2627
under the best configuration ($+0.1099$), comparable to the overall accuracy
increase ($+0.1260$). The concentration is real, but its cause lies in the type
composition of the evaluation sample rather than in per-type learning difficulty
--- diagnosis accounts for just 23 of all 6,719 PathVQA test items (0.34\%) ---
which is a dataset-level constraint rather than a defect of this evaluation
design.

\paragraph{What automated search reached, and the limits of autonomous agents.}
Phase~2's fixed-axis ablation and Phase~3's free search arrived independently at
the same region (\texttt{lora\_targets=full}, rank in 32--64), which raises
confidence in the Phase~2 conclusion and partially compensates for the ablation's
inability to validate the three optima in combination. There was, however, a
clear boundary to the level of automation: automation of the search algorithm
surpassed manual configuration by a wide margin (0.4490 vs.\ 0.3776), yet the LLM
agent's autonomous judgement fell significantly short of that algorithm. The
cause was not an absence of predictive ability --- the trial-level mean quality of
the agent's proposals was in fact the highest (0.3980) --- but a search-behaviour
problem, and that problem proved correctable: removing the configuration
inconsistencies restored unique combinations to 18.90 / 20 and eliminated the
significant deficit against Optuna. Restored diversity did not, however,
translate into performance; the run-level mean of 0.4292 still trailed Optuna
(0.4490), and only 2 of 10 repetitions reached the upper band. The remaining gap
therefore resembles a problem of judgement --- identifying promising regions and
allocating budget toward them --- rather than the procedural defect of repeating
the same configuration. Within the scope of this study, automation of search is
mature and the agent's procedural defects proved correctable, but whether
autonomous research judgement can replace existing optimization algorithms remains
unconfirmed.

\paragraph{The unit of aggregation changes the conclusion.}
On two separate occasions the choice of aggregation unit reversed the conclusion
itself. Pooling the four models, the mixed-effects model reported the fine-tuning
effect as not significant ($p = .3629$) --- not because there was no effect but
because effects in opposite directions cancelled in the pooled mean; decomposed
by model, three of the four show significant effects, two of them negative.
Treating trials as independent observations would likewise have reversed RQ3: the
trial-level mean is higher for Autoresearch (0.3980) than for Optuna (0.3905),
yet the run-level best is significantly higher for Optuna (0.4490) than for
Autoresearch (0.4184). The observation that re-running an identical configuration
produced \texttt{val\_accuracy} between 0.388 and 0.444 --- a range of about
0.056, larger than the mean difference between strategies detected here
($0.4490 - 0.4184 = 0.031$) --- makes the point quantitative: judging the relative
merit of strategies from a single trial result is impossible in principle, and a
design using 10 independent repeats as the unit of testing is the minimum
requirement given the scale of this noise.

\paragraph{Indirect comparison with prior work.}
We did not reproduce medical-specialized VLMs in our environment. LLaVA-Med
\citep{li2023llavamed} does report figures on the standard test splits of the same
three datasets, permitting an indirect comparison on closed-ended accuracy only
(Appendix~\ref{app:tables}, Table~\ref{tab:appllavamed}); open-ended figures are
not comparable, because LLaVA-Med scores open items by the proportion of
generated responses containing the gold token while we use a BERTScore F1
threshold; the former is structurally a far more lenient measure, so placing the
two figures side by side would itself be misleading. On SLAKE closed accuracy a 2B
model fine-tuned only with QLoRA reached practically the same level as the 7B
medical-specialized model (85.26 vs.\ 85.34), and since all of these greatly
exceed general-purpose LLaVA without domain training (63.22), this suggests that
for tasks at the level of SLAKE the advantage of large-scale domain pretraining
can be substantially offset by QLoRA adaptation alone. On PathVQA (83.12 vs.\
91.21) and VQA-RAD (72.91 vs.\ 84.19) a gap of 8--11\,\%p remains, so that
advantage remains clearly visible in proportion to the specialization of the task.
The comparison carries three constraints: our figures are means over three seeds
against a single-run report, the training budgets differ greatly, and the figures
are quoted from a published paper rather than reproduced with identical code in
an identical environment.

\paragraph{Limitations.}
Beyond those already stated, five constrain the results most directly. Phase~3
fixed \texttt{max\_steps}~$=200$ but not the number of samples actually seen:
because the search space includes \texttt{batch\_size} and
\texttt{grad\_accum\_steps}, effective training volume ranged from 800 to 12,800
samples, a factor of about 16, so the strategy comparison mixes training volume
with hyperparameter quality --- though the best configurations of Optuna and
Autoresearch were both batch 4 $\times$ grad\_accum 16 (12,800 samples), so the
comparison central to RQ3 is not explained by this confound. The agent was
non-deterministic: we originally stated that temperature was fixed at 0, and
post-hoc verification established that it was not; the implementation applies a
schedule and an exhaustive tally of 432 API calls shows the temperature actually
used ranged from 1.0 down to 0.66, with the 0 recorded in the results table being
an unpopulated default. Statistical power is limited ($n = 9$ in Phase~2, $n = 10$
per strategy in Phase~3), leaving effect-size intervals wide, and no
multiple-comparison correction was applied across the pipeline --- Bonferroni was
applied only to the Phase~1 McNemar post-hoc tests --- though the principal
conclusions rest on triple verification or non-overlapping confidence intervals
rather than single $p$-values. The ablations were conducted under the single
condition of PathVQA and Qwen3-VL-2B, and the planned supplementary rank
verification on SLAKE could not be carried out because of GPU time constraints.
Finally, feasibility on 16\,GB-class hardware was inferred from measured peak
VRAM (at most 14,373\,MB) rather than verified by execution on such a card; in a
real 16\,GB environment usable VRAM is smaller than nominal capacity and
fragmentation behaves differently, so an out-of-memory risk cannot be excluded
for Gemma4-E2B in particular, whose headroom is under 2\,GB, while the other
three models (4,015--7,943\,MB) have ample margin.
